\chapter{Introduction}
\pagenumbering{arabic} \setcounter{page}{1}

This project presents PropertyIQ: a web-based application designed to offer unambiguous insights into properties available on the market through the use of ML models, and it has been uploaded and can be found on GitHub \parencite{Kirilov2026}. This project introduces an \textbf{innovative methodology for comprehending property values} by leveraging known attributes that characterise and define a property while incorporating external factors such as the quality of the location.              

PropertyIQ is constructed using two complementary machine learning models: the first model predicts the value of a property based on the property features and facilitates a comparison with existing properties on the market, while the second model analyses the resulting valuation and generates human-readable explanations of the predicted price. 

Both models were trained on a dataset generated by scraping Zoopla \parencite{Data2024}, which contains real-world values from existing properties with detailed features.

This project aligns with Goal 9 of the Sustainable Development Goals of the United Nations: \textit{build resilient infrastructure, promote inclusive and sustainable industrialisation, and foster innovation} \parencite{UnitedNations2011}. PropertyIQ advances the \textbf{inclusive industrialisation} pillar by democratising access to property valuation, enabling users to obtain a credible estimate of a property's market value regardless of prior appraisal expertise and thereby reducing the knowledge required that has historically disadvantaged first-time buyers, small landlords, and less affluent market participants. It further contributes to the \textbf{innovation} pillar by adopting a system that, although well-established in the academic machine-learning literature, remains insufficiently exploited in the property market: model prediction explainability. Rather than returning a single opaque estimate, the system decomposes the prediction into the contribution of individual property features, allowing users to interrogate the given valuation that has been produced; as discussed in Section~\ref{se:GeneralBenfits} \nameref{se:GeneralBenfits}. This transparency is of substantial practical benefit to non-expert users, who require not just a figure but a justification on which to base a significant financial decision.


\section{Motivation and Objectives}

This project seeks to address the existing gap between the application of Artificial Intelligence (AI) systems that solely assess a property's value and the necessity to comprehend the rationale behind such valuations. PropertyIQ aims to rectify this deficiency by integrating predictive modelling with XAI techniques, enabling users to engage with both the valuation process and its underlying justification. 

The following objectives for this project have been identified.

\begin{enumerate}

  \item \textbf{Achieve accurate property valuations} using ML, by training a regression model on real-world UK property data to produce reliable value estimates based on a defined set of property features.
  
  \item \textbf{Ensure transparent and comprehensible explanations} by employing a machine learning technique to decompose each valuation into the individual contributions of each feature, which are then translated into plain English that is accessible to non-technical users.
  
  \item \textbf{Provide a comprehensive understanding of valuations} within the market; it is essential to present each predicted value in conjunction with comparable properties from the dataset. This will also allow users to assess the estimate by comparing it with similar listings in the same geographical area rather than evaluating it in isolation.
  
  \item \textbf{Facilitate informed decision-making} in the property market by elucidating the factors underlying property valuations and encourage individuals to critically evaluate property estimates rather than accepting them uncritically, thus mitigating the risk of overconfidence, which may result in excessive expenditure over time.
  
  \item \textbf{Develop an accessible web-based interface} that is suitable for prospective buyers, small-scale investors, and property researchers without requiring technical expertise in machine learning or property investment.
  
  \item \textbf{Establish a scalable and maintainable framework} in which the underlying dataset, model, and explanation pipeline can be extended on new data. This approach will facilitate future development beyond the scope of this project.
  
\end{enumerate}


